% СПО ЛКС - безопасность сетей -- классификация
% Пынькин Д.А. (с) 2012
\documentclass[ignorenonframetext, hyperref={pdftex, unicode}]{beamer}

\usepackage[russian]{babel}
\usepackage[utf8]{inputenc}
\usepackage[T1]{fontenc}

\usepackage{beamerthemesplit}

%\usetheme{Pittsburgh}
%\usecolortheme{dolphin}

\usetheme{Copenhagen}
\usecolortheme{wolverine}


\usepackage{ulem}
%\usepackage{html}

\usepackage{verbatim}

\usepackage{ucs}
\usepackage{listings}
\lstloadlanguages{C, make, bash}

\lstset{escapechar=`,
	extendedchars=false,
	language=C, 
	tabsize=2, 
	columns=fullflexible, 
%	basicstyle=\scriptsize,
	keywordstyle=\color{blue}, 
	commentstyle=\itshape\color{brown},
%	identifierstyle=\ttfamily, 
	stringstyle=\mdseries\color{green}, 
	showstringspaces=false, 
	numbers=left, 
	numberstyle=\tiny, 
	breaklines=true, 
	inputencoding=utf8x,
	keepspaces=true,
	morekeywords={u\_short, u\_char, u\_long, in\_addr}
	}

\definecolor{darkgreen}{cmyk}{0.7, 0, 1, 0.5}

\lstdefinelanguage{diff}
{
    morekeywords={+, -},
    sensitive=false,
    morecomment=[l]{//},
    morecomment=[s]{/*}{*/},
    morecomment=[l][\color{darkgreen}]{+},
    morecomment=[l][\color{red}]{-},
    morestring=[b]",
}

\title[СПОЛКС (http://goo.gl/32cTB)]{Системное программное обеспечение\\локальных компьютерных сетей}
\author{Денис Пынькин}
\date{2011 -- 2012}
%\institution[БГУИР]{Белорусский государственный университет информатики и радиоэлектроники}
%\logo{\includegraphics[width=1cm]{logo-kafEVM.png}}


\subtitle{Обзор и классификация\\сетевых угроз безопасности}

\begin{document}

\mode<all>{\begin{frame}
\titlepage
\begin{center}
e-mail: denis.pynkin@bsuir.by\\
\end{center}
\begin{center}
{\bfseries http://goo.gl/32cTB}

{\tiny СЧАСТЬЕ ДЛЯ ВСЕХ, ДАРОМ, И ПУСТЬ НИКТО НЕ УЙДЕТ ОБИЖЕННЫЙ!\\
(c)Стругацкие, Пикник на обочине}
\end{center}
\end{frame}
}

%
% Далее начинается сама презентация
%

\begin{frame}{Введение}
Рассматриваются удаленные атаки на инфраструктуру и протоколы распределенных вычислительных систем.
\end{frame}


\begin{frame}{Терминология}
	\begin{block}{Субъект атаки (источник атаки)}
		это атакующая программа или оператор,  непосредственно осуществляющие воздействие
	\end{block}
	\begin{block}{Типовая угроза безопасности}
		удаленное информационное разрушающее воздействие,  программно осуществляемое по каналам связи и характерное для любой распределенной ВС.
	\end{block}
	\begin{block}{Типовая удаленная атака}
		реализация типовой угрозы безопасности РВС.
	\end{block}
\end{frame}


\begin{frame}{Подмножества -- по объекту атаки}

	\begin{itemize}
		\item удаленные атаки на инфраструктуру и протоколы сети;
		\item удаленные атаки на телекоммуникационные службы или серверы предоставления удаленного сервиса.
	\end{itemize}
\end{frame}


\begin{frame}{Основные причины}
Благодаря которым возможна реализация удаленных угроз:
\begin{itemize}
	\item использование широковещательной среды передачи (например,  Ethernet); \pause
	\item применение нестойких алгоритмов идентификации удаленных субъектов и объектов РВС; \pause
	\item использование протоколов динамического изменения маршрутизации с нестойкими алгоритмами идентификации; \pause
	\item применение алгоритмов удаленного поиска с использованием широковещательных и направленных поисковых запросов; \pause
	\item возможность анонимного захвата одним субъектом РВС множества физических или логических каналов связи.
\end{itemize}
\end{frame}


\section{Классификация}
\begin{frame}{}
	\huge
	\center
	Общая классификация угроз сетевой безопасности

	\bigskip
	\normalsize
	По материалам книги: ''АТАКА НА INTERNET''

	И.Д. Медведовский, П.В. Семьянов, Д.Г. Леонов

\end{frame}


\begin{frame}{1. По характеру воздействия}
	\begin{enumerate}
		\item пассивное -- не оказывает непосредственного влияния на работу системы,  но способно нарушать ее политику безопасности;
		\item активное -- воздействие,  оказывающее непосредственное влияние на работу системы и нарушающее принятую в ней политику.
	\end{enumerate}
\end{frame}


\begin{frame}{2. По цели воздействия}
	\begin{enumerate}
		\item нарушение конфиденциальности информации либо ресурсов системы;
		\item нарушение целостности информации;
		\item нарушение работоспособности (доступности) системы. 
	\end{enumerate}
	\bigskip
	Цель большинства атак -- получить несанкционированный доступ к информации:
	\begin{itemize}
		\item перехват -- получение информации без возможности ее искажения;
		\item искажение информации -- полный контроль над информационным потоком между объектами системы или возможность передачи сообщений от имени другого объекта.
	\end{itemize}
\end{frame}


\begin{frame}{3. По условию начала осуществления воздействия}
	\begin{enumerate}
		\item атака после запроса от атакуемого объекта;
		\item атака после наступления ожидаемого события на атакуемом объекте;
		\item безусловная атака. 
	\end{enumerate}
\end{frame}


\begin{frame}{4. По наличию обратной связи с атакуемым объектом}
	\begin{enumerate}
		\item с обратной связью;
		\item без обратной связи (однонаправленная атака).
	\end{enumerate}
\end{frame}


\begin{frame}{5. По расположению субъекта атаки относительно атакуемого объекта}
	\begin{enumerate}
		\item внутрисегментное;
		\item межсегментное.
	\end{enumerate}

	Подсеть --  логическое объединение, совокупность хостов,  являющихся частью глобальной сети,  для которых маршрутизатором выделен одинаковый номер. Хосты внутри одной подсети могут взаимодействовать непосредственно,  минуя маршрутизатор.

	Сегмент сети -- физическое объединение хостов.

\end{frame}


\begin{frame}{6. По уровню эталонной модели ISO/OSI}
	\begin{enumerate}
		\item физический;
		\item канальный;
		\item сетевой;
		\item транспортный; 
		\item сеансовый; 
		\item представительный; 
		\item прикладной. 
	\end{enumerate}
\end{frame}


\begin{frame}{Классификация типовых удаленных атак}
	\tiny
%	\fontsize{5}{7pt}\selectfont
%	\addtolength{\textwidth}{10cm}
	\begin{table}[ht]
	\centering
	\begin{tabular}[c]{|p{0.5cm}||*2{p{0.1cm}|}|*3{p{0.1cm}|}|*3{p{0.1cm}|}|*2{p{0.1cm}|}|*2{p{0.1cm}|}|*7{p{0.1cm}|}}
%\begin{tabular}[c]{|p{0.5cm}||*2{c|}|*3{c|}|*3{c|}|*2{c|}|*2{c|}|*3{c|}|*7{c|}}

%Типовая удаленная атака & Характер воздействия & Цель воздействия &Условие начала осуществления воздействия & Наличие обратной связи с атакуемым объектом & Расположение субъекта атаки относительно атакуемого объекта & Уровень модели OSI \\
	\hline
 & 1.1 & 1.2 & 2.1& 2.2 & 2.3 & 3.1 & 3.2 & 3.3 & 4.1 & 4.2 & 5.1 & 5.2 & 6.1 & 6.2 & 6.3 & 6.4 & 6.5 & 6.6 & 6.7 \\
	\hline
	\hline
1 &+&-&+&-&-&-&-&+&-&+&+&-&-&+&-&-&-&-&-\\
	\hline
2 &-&+&+&+&-&-&+&-&+&+&+&+&-&-&+&+&-&-&-\\
\hline
3 &-&+&+&+&+&-&-&+&+&+&+&+&-&-&+&-&-&-&-\\
\hline
4 &-&+&+&+&-&+&-&+&+&-&+&+&-&+&+&+&-&-&-\\
\hline
5 &-&+&-&-&+&-&-&+&-&+&+&+&-&+&+&+&+&+&+\\
\hline
	\end{tabular}
	\end{table}

	\normalsize
	\begin{enumerate}
		\item Анализ сетевого трафика
		\item Подмена доверенного объекта РВС
		\item Внедрение в РВС ложного объекта путем навязывания ложного маршрута
		\item Внедрение в РВС ложного объекта путем использования недостатков алгоритмов удаленного поиска
		\item Отказ в обслуживании
	\end{enumerate}
\end{frame}

\section{Классификация web-угроз}
\begin{frame}{}

	\begin{block}{Web Application Security Consortium}

	Данная классификация представляет собой совместную попытку собрать воедино и организовать угрозы безопасности Web серверов. Члены Web Application Security Consortium создали данный проект для разработки и популяризации стандартной терминологии  описания этих проблем.  Это даст возможность разработчикам приложений,  специалистам в области безопасности,  производителям программных продуктов и аудиторам  использовать единый язык для взаимодействия.
	\end{block}
\end{frame}


\begin{frame}{Цели}
	\begin{itemize}
		\item Определить все известные классы атак на Web-приложения.
		\item Согласовать названия для каждого из классов.
		\item Разработать структурированный подход к классификации  атак.
		\item Разработать документацию,  содержащую общее описание каждого из классов.
	\end{itemize}
\end{frame}


\begin{frame}{Аутентификация (Authentication)}
	\begin{itemize}
		\item Подбор (Brute Force) 
		\item Недостаточная аутентификация(Insufficient Authentication)
		\item Небезопасное восстановление паролей (Weak Password Recovery Validation)
	\end{itemize}

\end{frame}


\begin{frame}{Авторизация (Authorization)}
	\begin{itemize}
		\item Предсказуемое значение идентификатора сессии (Credential/Session Prediction)
		\item Недостаточная авторизация (Insufficient Authorization)
		\item Отсутствие таймаута сессии (Insufficient Session Expiration)
		\item Фиксация сессии (Session Fixation)
	\end{itemize}

\end{frame}


\begin{frame}{Атаки на клиентов (Client-side Attacks)}
	\begin{itemize}
		\item Подмена содержимого (Content Spoofing)
		\item Межсайтовое выполнение сценариев (Cross-site Scripting,  XSS)
		\item Расщепление HTTP-запроса (HTTP Response Splitting)
	\end{itemize}

\end{frame}


\begin{frame}{Выполнение кода (Command Execution)}
	\begin{itemize}
		\item Переполнение буфера (Buffer Overflow)
		\item Атака на функции форматирования строк (Format String Attack)
		\item Внедрение операторов LDAP (LDAP Injection)
		\item Выполнение команд ОС (OS Commanding)
		\item Внедрение операторов SQL (SQL Injection)
		\item Внедрение  серверных сценариев (SSI Injection)
		\item Внедрение операторов XPath (XPath Injection)
	\end{itemize}

\end{frame}


\begin{frame}{Разглашение информации (Information Disclosure)}
	\begin{itemize}
		\item Индексирование директорий (Directory Indexing)
		\item Идентификация приложений (Web Server/Application Fingerprinting)
		\item Утечка информации (Information Leakage)
		\item Обратный путь в директориях (Path Traversal)
		\item Предсказуемое расположение ресурсов (Predictable Resource Location)
	\end{itemize}

\end{frame}


\begin{frame}{Логические атаки (Logical Attacks)}
	\begin{itemize}
		\item Злоупотребление функциональными возможностями (Abuse of Functionality)
		\item Отказ в обслуживании (Denial of Service)
		\item Недостаточное противодействие автоматизации (Insufficient Anti-automation)
		\item Недостаточная проверка процесса (Insufficient Process Validation)
	\end{itemize}

\end{frame}
\mode<all>{\begin{frame}{}
\Huge
\begin{center}
	Спасибо за внимание!
	\bigskip
	Вопросы?
\end{center}
\normalsize
\end{frame}
}

\end{document}
%Конец файла
